\chapter{Płytka pośrednicząca w komunikacji}
    Aby usprawnić pobieranie danych przy pomocy kamizelki BRV utworzona została płytka pośrednia, która ma za zadanie odbierać dane przekazywane szeregowo i wysyłać je dalej przy pomocy technologii Bluetooth. W tym rozdziale omówiona zostanie budowa płytki, oraz struktura zaimplementowanego oprogramowania.
    W tym momencie płytka pośrednia udostępnia następujące usługi:
    \begin{itemize}
        \item {\textbf{BRICS BLE} - serwer Bluetooth LE w trybie Peripheral}
        \item {\textbf{BRICS Serial} - wielowątkowa aplikacja pobierająca szeregowo dane z kamizelki BRV}
        \item {\textbf{BRICS.service}}
    \end{itemize}
    Usługi te są dostępne razem z podłączoną kamizelką BRV i mają na celu jednynie usprawnienie przekazywania danych i ograniczenie utraty danych podczas pomiarów.

    \section {Architektura płytki}
        Płytka została zrealizowana przy pomocy mikroprocesora Raspberry Pi Zero W, na którym zainstalowany został system operacyjny Raspbian 32-bit. Posiada 64 GB pamięć w postaci karty micro SD.
        \\ \\


    \section {Wielowątkowa aplikacja pośrednicząca}
        Do obsługi współbieżnego przetwarzania danych zaimplementowana została wielowątkowa aplikacja. \\
        Uruchomienie aplikacji powoduje uruchomienie trzech wątków:
        \begin{itemize}
            \item{\textbf{BRVBluetoothServer} - wątek obsługujący działanie serwera Bluetooth, którego specyfikacja zostanie omówiona w dalszej części dokumentu.}
            \item{\textbf{BRVSerialConnection} - wątek obsługujący szeregowe pobieranie pakietów danych z kamizelki BRV.}
            \item{\textbf{BRVLedDisplay} - wątek obsługujący sygnalizowanie przy pomocy dołączonych diod o aktywności i stanie pracy kolejnych komponentów.}
        \end{itemize}
        Komunikacja pomiędzy wątkami została zrealizowana przy pomocy bibliotek:
        \begin{itemize}
            \item{\textbf{queue} - biblioteka udostępniająca kolejkę FIFO z wbudowanym mechanizmem blokowania}
            \item{\textbf{asyncio} - biblioteka umożliwiająca asynchroniczne uruchamianie kodu}
            \item{\textbf{threading} - biblioteka umożliwiająca tworzenie i kontrolowanie wątków}
        \end{itemize}
        Dane otrzymane na porcie szeregowym są umieszczane w tej kolejce i zdejmowane przez wątek obsługujący serwer Bluetooth. W przypadku, w którym aplikacja napotka błąd lub nieobsłużony wyjątek wątki są niszczone i aplikacja jest uruchamiana ponownie.
    \section {Serwer Bluetooth}
        W celu utworzenia serwera BLE wykorzystane zostały biblioteki:
        \begin{itemize}
            \item{\textbf{bluez} - podstawowa biblioteka służaca do tworzenia połączeń wykorzystujących technologię Bluetooth}
            \item{\textbf{bluez-peripheral} - bliblioteka służąca do tworzenia serwerów BLE wykorzystujących bibliotekę Bluez, oraz GATT API}
            \item{\textbf{dbus} - biblioteka służąca do komunikacji z wykorzystaniem głównej szyny danych}
        \end{itemize}
        Utworzony serwer składa się z dwóch klas, jedna z nich stanowi deklaracje serwisu udostępnianego przez serwer Bluetooth, a druga obsługę wątka uruchamiającego serwer Bluetooth.
        \\ \\
        Serwer działa w 2 następujących po sobie fazach:
        \begin{enumerate}
            \item{\textbf{Uruchamianie} - podczas tej fazy serwer musi zadeklarować swoje serwisy, charakterystyki, adapter, agent i rozgłaszanie. Dane o serwisie i charakterystyce dostępnej na płytce definiowane są w klasie \textbf{BRVDataService}, a sposób rozgłaszania, agent oraz adapter definiowane są podczas 
            inicjalizacji serwera w klasie \textbf{BRVBluetoothServer}. Ponadto, podczas inicjalizacji wątek dostaje od systemu szynę danych d-bus, na której wykonywana będzie konieczna komunikacja ze sterownikami Bluetooth i kolejno inicjalizowane są kolejne elementy serwera.}
            \item{\textbf{Działanie} - podczas tej fazy kolejno zdejmowane są kolejne pakiety z kolejki FIFO, oraz rozdzielane na mniejsze aby umożliwić ich transport przy pomocy Bluetooth. Aby wysłać pakiet na dostępnej charakterystyce serwisu uruchamiana jest metoda \textbf{update\_BRV\_value}, która powoduje wygenerowanie
            powiadomienia. Wszystkie urządzenia, które są podłączone pod system powiadomień tej charakterystyki otrzymują informacje o jej zmianie i mogą odczytać wartość}
        \end{enumerate}
            \subsection{Działanie serwera Bluetooth LE}
    Serwer Bluetooth LE (Low Energy) składa się z kilku znaczących komponentów. Głównym z nich jest serwis, który zgodnie z protokołem GATT serwis zawiera w sobie charakterystykę. Jest ona ustawiona w tryb Indicate, co umożliwia otrzmywanie komunikatów ACK po wysłaniu pakietów.
    \section {Sygnalizowanie stanu płytki GPIO}
        W ramach aplikacji sygnalizującej stan pracy płytki obsługiwane jest działanie 3 diod: sygnalizującej prace płytki, sygnalizującej odbieranie danych przez port szeregowy, oraz sygnalizującej działanie serwera Bluetooth. Sterowanie działaniem diod odbywa się przy pomocy map bitowych, których stan jest regularnie brany pod uwagę podczas działania programu.
        \\ \\
        Możliwe stany diod:
        \begin{itemize}
            \item{dioda sygnalizująca działanie płytki świeci się stałym światłem - płytka działa}
            \item{dioda sygnalizująca działanie płytki mruga - działa aplikacja obsługująca przekazywanie danych}
            \item{dioda sygnalizująca działanie serwera mruga - działa serwer Bluetooth i dane są aktywnie przekazywane}
            \item{dioda sygnalizująca działanie portu szeregowego mruga - pobierane są dane z portu szeregowego}
        \end{itemize}
    \section {Przykładowy scenariusz użycia}
        Płytka może zostać użyta w następujący sposób:
        \begin{enumerate}
            \item{Podłaczenie płytki do zasilania}
            \item{Uruchomienie programu odbierającego dane}
            \item{Zapaliła się dioda sygnalizująca działanie serwera Bluetooth i w programie odbierającym dane widać podłączenie do serwera}
            \item{Podłaczenie płytki do kamizelki BRV i rozpoczęcie pomiaru}
            \item{Zakończenie pomiaru i odpięcie kamizelki BRV od zasilania}
            \item{Odłaczenie płytki od prądu}
        \end{enumerate}
